Тема дипломной работы: <<Разработка программной системы фиксации городских событий>>.

В рамках преддипломной практики необходимо систематизировать и обобщить результаты
двух этапов разработки программной системы фиксации городских событий для жителей
города Кирова, выполненных в ходе курсовых работ.

\subsection{Этап 1. Серверная часть и веб-панель диспетчера}

На первом этапе (курсовая работа, 1-й семестр) были поставлены и решены следующие задачи:
\begin{enumerate}
    \item исследование текущего процесса подачи и обработки обращений граждан по вопросам
          городской инфраструктуры города Кирова, выявление его ключевых проблем~---
          отсутствие единого канала обращений, отсутствие обратной связи о статусе
          и невозможность визуализации проблем на карте;
    \item проведение сравнительного анализа существующих аналогов (ПОС <<Госуслуги.
          Решаем вместе>>, <<Наш город>>, <<Добродел>>, <<РосЯма>>, FixMyStreet)
          и формулирование требований к системе;
    \item проектирование трёхзвенной клиент-серверной структуры системы
          (клиентские приложения~--- сервер приложений~--- база данных) в виде модульного
          монолита с использованием чистой структуры;
    \item проектирование и реализация реляционной базы данных PostgreSQL
          с ограничениями целостности, индексами и поддержкой геопространственных
          запросов через расширение PostGIS;
    \item реализация серверного программного интерфейса на платформе ASP.NET Core,
          включая алгоритм поиска дубликатов обращений в радиусе 50~м и маршрутизацию
          обращений к ответственным службам;
    \item проектирование и реализация прототипа диспетчерской веб-панели на React с
          библиотекой Яндекс.Карты, включая дашборд со сводной статистикой,
          интерактивную карту обращений с фильтрацией, очередь обращений
          с вкладками по статусам, карточку обращения с историей статусов,
          справочники категорий проблем и городских служб;
    \item тестирование серверной части (модульные и интеграционные тесты на xUnit,
          тестирование API через Swagger UI) и функциональное тестирование
          веб-панели диспетчера.
\end{enumerate}

Компоненты системы, реализованные на данном этапе:
\begin{itemize}
    \item реляционная база данных PostgreSQL с расширением PostGIS для хранения
          обращений, категорий проблем, справочника служб, истории статусов
          и медиафайлов обращений;
    \item серверный программный интерфейс (REST API) на платформе ASP.NET Core
          с Entity Framework Core, обеспечивающий взаимодействие клиентских приложений
          с базой данных, аутентификацию, авторизацию, валидацию данных
          и геопространственные запросы;
    \item диспетчерская веб-панель~--- автоматизированное рабочее место оператора
          на React и TypeScript с библиотекой Яндекс.Карты для интерактивной карты, сборщиком
          Vite и системой управления состоянием на основе React Context.
\end{itemize}

\subsection{Этап 2. Мобильное приложение для жителей}

На втором этапе (курсовая работа, 2-й семестр) были поставлены и решены следующие задачи:
\begin{enumerate}
    \item проведение анализа предметной области с акцентом на мобильный канал подачи
          обращений и формулирование требований к мобильному приложению с учётом
          разработанной ранее серверной части;
    \item моделирование текущего и планируемого процессов фиксации городских проблем
          в нотациях IDEF0 и IDEF3 с выявлением места мобильного приложения
          в общей структуре системы;
    \item выполнение обзора аналогов с акцентом на мобильные приложения~---
          <<Госуслуги. Решаем вместе>>, <<Наш город>>, <<Добродел>>,
          <<Наш Санкт-Петербург>>, FixMyStreet~--- с анализом кроссплатформенности,
          GPS-функций, push-уведомлений и механизмов дедупликации;
    \item разработка технического задания на мобильное приложение: роли пользователей,
          функциональные и нефункциональные требования, требования к составу экранов;
    \item проектирование структуры мобильного приложения на React Native с Expo,
          включая структуру экранов, навигацию, клиент API и сервисы устройства;
    \item проектирование пользовательского интерфейса~--- разработка макетов
          экранов авторизации, карты, создания обращения, списка обращений,
          карточки обращения, уведомлений и профиля;
    \item реализация кроссплатформенного мобильного приложения на React Native
          с TypeScript и инструментарием Expo, включая интеграцию с камерой
          и геолокацией устройства;
    \item проведение модульного тестирования (20 тестовых файлов на Vitest),
          тестирование API и функциональное тестирование мобильного приложения.
\end{enumerate}

Компоненты системы, реализованные на данном этапе:
\begin{itemize}
    \item кроссплатформенное мобильное приложение для платформ Android и iOS на базе
          React Native с TypeScript и инструментарием Expo, предназначенное
          для подачи обращений жителями;
    \item интеграция мобильного приложения с существующим серверным программным
          интерфейсом: аутентификация на основе JWT, отправка обращений с медиафайлами,
          получение статусов и уведомлений;
    \item доработка серверной части для поддержки функций мобильного приложения~---
          эндпоинт поиска дубликатов по координатам и категории, подтверждение
          обращений другими пользователями, загрузка медиафайлов.
\end{itemize}

\subsection{Обобщение задания}

Таким образом, в ходе преддипломной практики необходимо обобщить результаты проектирования,
реализации и тестирования полной программной системы фиксации городских событий, включающей
три основных компонента: серверную часть с базой данных, диспетчерскую веб-панель и мобильное
приложение для жителей. На рисунке~\ref{fig:architecture-full} представлена диаграмма
компонентов, отображающая структуру программной системы и взаимосвязи между её составными
частями.

\begin{figure}[H]
    \centering
    \includegraphics[width=\textwidth]{architecture-full.png}
    \caption{Диаграмма компонентов программной системы фиксации городских событий}
    \label{fig:architecture-full}
\end{figure}

Итогом практики является подготовка материалов для пояснительной записки к выпускной
квалификационной работе.
